\begin{hyperbolicbox}{Summary}
    In this poster, we explored the action $\operatorname{PSL}(2,\Z)$ on $\H$. We observed that the action of $\operatorname{PSL}(2,\Z)$ tessellates $\H$ into copies of a fundamental domain, where any two triangles can be associated by a M\"obius map in $\operatorname{PSL}(2,\Z)$.

    Then we studied the lattice structure in $\C$: any lattice is generated by independent $p,q\in\C$. While homothety provides a notion of similarity among lattices, it actually generates an equivalence relation, based on which one can classify all complex lattices.

    We finally discovered that $\Lambda_\tau$ is homothetic to $\Lambda_{\tau'}$ if and only if $\tau'$ can be transformed from $\tau$ via a M\"obius map in $\operatorname{PSL}(2,\Z)$. It follows that the orbit space $\H/\operatorname{PSL}(2,\Z)$ and the set of lattices up to homothety have the same structure and are connected by a bijection $\tau\mapsto\Lambda_\tau$. The $j$-invariant, as a parameterization of equivalence classes, gives us an algebraic view of this connection.
\end{hyperbolicbox}